\section{问题分析}

\subsection{机理分析：三个物理关系与三层成本}

微网调度的物理内核可归结为三个关系与三层成本\cite{ieee_tou,thu_ess}。

\textbf{（1）逐时段功率平衡（允许弃光）}。任一 10 分钟时段 $t$ 内，微网供给必须不低于小区负载：
\begin{equation}
PV_t + P_t + D_t - C_t \ \ge\ L_t ,
\label{eq:balance_general}
\end{equation}
其中 $P_t$ 为向外部电网购电量，$C_t,D_t$ 分别为储能充电量与放电量，$PV_t$ 为光伏发电量，$L_t$ 为负载电量。
取“$\ge$”而非“$=$”是因为实测数据中存在光伏富余功率超过储能最大充电功率（5000 kW，即
每时段 833.33 kWh）的情形，此时必须允许\textbf{弃光}，否则模型无可行解（该现象在全年出现于
多个晴好低负荷日）。

\textbf{（2）储能荷电状态（SOC）递推}。设储电量 $E_t$（kWh），充放电效率按单程 90\% 计：
\begin{equation}
E_{t+1} = E_t + \eta_c C_t - \frac{D_t}{\eta_d}, \qquad \eta_c=\eta_d=0.9 ,
\label{eq:soc}
\end{equation}
即充电 1 kWh 仅存入 0.9 kWh、放出 1 kWh 需消耗 $1/0.9$ kWh 储电量，往返效率 81\%。

\textbf{（3）物理边界}。
\begin{equation}
1200 \le E_t \le 10800,\qquad 0\le C_t \le \bar C,\qquad 0\le D_t \le \bar D,\qquad \bar C=\bar D=5000\times\frac{10}{60}=833.33 .
\label{eq:bounds}
\end{equation}

\textbf{（4）三层成本结构}：①正常购电按交易时刻电价（问题一、二、三中为附件1 固定电价曲线，
问题四中为附件4 实时电价）；②实际供电低于负载的缺口按 \textbf{5 倍}电价紧急购电（偏差电量惩罚结算机制见文献\cite{cn111260354}）；
③问题三中计划量与调整量的偏差按 \textbf{0.5 倍}（计划超出、退费 50\%）与
\textbf{1.5 倍}（调整超出）分段计价。

\subsection{四问的信息结构与决策结构差异}

四问的物理模型相同，差别在于\textbf{计划时可获得的信息}与\textbf{结算时的计价规则}，见\cref{tab:info}。

\begin{table}[htbp]
\centering
\caption{四问的信息集与结算规则对比}
\label{tab:info}
\small
\begin{tabularx}{\textwidth}{@{}p{1.2cm}p{2.5cm}p{3.3cm}p{3.6cm}X@{}}
\toprule
问题 & 电价口径 & 0:00 计划的信息集 & 计划/调整/紧急计价 & 求解窗口 \\
\midrule
一 & 附件1 固定 & 附件1 负载与光伏预测（整日已知） & 计划购电价 & 单日 144 时段 \\
二 & 附件1 固定 & \textbf{不得使用当日实测}；用历史实测构造预测 & 计划购电 + 缺口 5 倍 & 2025.2.1--12.31（334 天） \\
三 & 附件1 固定 & 0:00 发布的光伏预报 + 历史负载（6/12/18 另有新预报可用） & 计划 + 调整（0.5/1.5）+ 缺口 5 倍 & 同上（0:00 计划 $+$ 按需在某一预报时刻调整一次） \\
四 & 附件4 波动 & 同问题三（光伏预报 + 历史负载），电价已知 & 同问题二、三 & 同上 \\
\bottomrule
\end{tabularx}
\end{table}

\begin{figure}[htbp]
\centering
\begin{tikzpicture}[
  node distance=6mm and 7mm,
  box/.style={rectangle, rounded corners=1.5pt, draw=black!70, fill=blue!6, align=center,
              inner sep=3pt, font=\scriptsize, minimum height=7mm, text width=24mm},
  plain/.style={rectangle, draw=black!55, fill=orange!8, align=center,
              inner sep=3pt, font=\scriptsize, minimum height=7mm, text width=24mm},
  arr/.style={-{Latex[length=2mm]}, draw=black!70}
]
\node[plain] (d1) {附件1/2/4\\10 分钟数据};
\node[plain, right=of d1] (d2) {附件3\\整点光伏预报};
\node[box, below=of d1] (pat) {规律检验\\周内/季节/噪声};
\node[box, below=of d2] (basis) {计划依据\\升级预测模型\\+ 报童裕度};
\node[box, below=of pat] (lp) {单日/剩余时段 LP\\$\text{min}\sum p_tP_t - vE_{T+1}$};
\node[box, below=of basis] (adj) {12:00 单次滚动调整\\（对照 6/18 时刻）\\偏差分段计价};
\node[box, below=of lp] (settle) {实时平衡控制\\先调计划充放\\再动用储能\\缺口 5 倍紧急购电};
\node[box, below=of adj] (out) {结果文件\\result1/2/3/4-2/4-3};
\draw[arr] (d1) -- (pat); \draw[arr] (d1) -- (basis);
\draw[arr] (d2) -- (basis); \draw[arr] (pat) -- (basis);
\draw[arr] (basis) -- (lp); \draw[arr] (basis) -- (adj);
\draw[arr] (lp) -- (settle); \draw[arr] (adj) -- (settle);
\draw[arr] (settle) -- (out); \draw[arr] (adj) -- (out);
\end{tikzpicture}
\caption{总体技术路线：数据规律 $\to$ 计划依据 $\to$ 线性规划 $\to$ 滚动调整 $\to$ 实时结算}
\label{fig:framework}
\end{figure}

\subsection{方法选型与理由}

\begin{enumerate}[label=(\arabic*), leftmargin=2.2em, itemsep=2pt]
\item \textbf{为什么用线性规划（LP）而非启发式算法}：本文模型的目标函数（购电费、紧急购电费、
偏差费用）与全部约束（功率平衡、SOC 递推、限幅）均为线性，且决策变量连续；
LP 可在毫秒级给出\textbf{全局最优解}。对照实验显示，启发式“贪心配对”在典型日比 LP 贵
7\,060.98 元（20.1\%），故启发式在本问题中只作为对照方法。
\item \textbf{为什么用报童安全裕度而非随机规划}：0:00 计划的信息不完美，买多按基准价浪费、
买少按 5 倍紧急购电，二者构成典型报童问题，理论最优服务水平 $c_u/(c_u+c_o)=0.8$；
因日内实时平衡控制已能吸收约一半的意外偏差，实测最优分位降为 0.7。相对地，用近 20 日实测情景构造的
两阶段随机规划（SAA）因“情景内储能自由度被重复利用”而严重低估尾部风险，
全年费用 39\,954\,010 元（比采用方案高 192.6\%，见第~\ref{sec:compare} 节）。
\item \textbf{为什么用滚动时域（MPC）调整}\cite{mpc,cn102184475}：预报随发布时刻临近而变准（同目标小时每隔 6 小时
预报更新幅度平均 273 kW），据此在每个预报时刻只对剩余时段重解 LP 并继承当前 SOC，
是标准滚动时域做法；该结构也天然对应题目给出的三段计价规则。
\item \textbf{为什么把“实时平衡控制”作为结算机制}：题面要求“微网提供的电能不可低于小区负载”，
即微网会主动用储能兜底，只有计划购电与储能都无法覆盖时才购买 5 倍紧急电；
本文把这一过程写成逐时段因果的控制规则（先调整计划充放电、再动用储能存量，
见式 \eqref{eq:clip}--\eqref{eq:socrec}），只用当前实测与当前储电量，因而\textbf{可实行}。
对照实验表明该机制比“按计划僵硬执行储能”节约 6.20\%（90.2 万元/年）。
\end{enumerate}
